
This work addressed whether data attribution method performance depends on model architecture, finding that it does---but the patterns differ from theoretical predictions.

TRAK exhibits architecture-invariance (<1\% AUC difference at all compute budgets). TracIn shows a directional BERT advantage (~3\% at low compute), though not statistically significant with the sample size used. EK-FAC does not show the predicted GPT-2 advantage (0.27\% difference, p=0.90).

The causal mechanism is validated: 98.82\% attention sparsity difference between architectures creates an $11\times$ Hessian eigenvalue difference, propagating through each method's approximation differently.

For practitioners choosing attribution methods across architectures, TRAK provides consistent performance. Architecture-agnostic attribution is achievable through random projection.

\subsubsection{Future Work}

\begin{itemize}
\item Scale testing with 1B+ parameter models
\item Task generalization to MNLI, SQuAD, and other benchmarks
\item Encoder-decoder architectures (T5)
\item Adequately-powered studies with 5+ seeds for P1/P2 confirmation
\end{itemize}
